\chapter{Introducción} \label{chap:introduccion}
El desarrollo de sistemas robóticos representa un campo interdisciplinario dentro de la ingeniería, que combina conocimientos de mecánica, electrónica, informática y control automático. En este reporte se documenta el diseño y la implementación de un robot con cuatro grados de libertad (4 GDL), abarcando los aspectos fundamentales de su modelado y operación mediante el análisis de la cinemática, la dinámica, el control y la implementación en ROS (Robot Operating System).

La cinemática del robot, tanto directa como inversa, es esencial para establecer la relación entre los movimientos articulares y la posición y orientación del efector final. A través del modelado cinemático se diseñan trayectorias deseadas en el espacio de trabajo. La dinámica, por su parte, permite entender las interacciones entre las fuerzas y movimientos, siendo fundamental para simular comportamientos reales y para el diseño de controladores robustos. En cuanto al control, se implementan estrategias que aseguren el seguimiento preciso de trayectorias, compensando perturbaciones y no linealidades del sistema. Finalmente, ROS se utiliza como plataforma de desarrollo y simulación, permitiendo la integración modular de componentes de software, visualización, y comunicación entre nodos del sistema robótico.

Este trabajo tiene como objetivo demostrar la aplicación práctica de estos conceptos teóricos en el desarrollo de un robot funcional, evidenciando la importancia de un enfoque sistemático e interdisciplinario en el área de la robótica moderna.
